\documentclass[10pt,twocolumn]{article}
\usepackage[a4paper,margin=1.35cm]{geometry}
\usepackage{graphicx,booktabs,amsmath,amssymb,microtype}
\usepackage{mathptmx}  % Times: ~8% denser than the default font
\usepackage[font=small,skip=2pt]{caption}
\usepackage[dvipsnames]{xcolor}
\usepackage{titlesec}
\titlespacing*{\section}{0pt}{4pt}{2pt}
\titlespacing*{\subsection}{0pt}{3pt}{1pt}
\titleformat{\section}{\normalfont\large\bfseries}{\thesection}{0.5em}{}
\titleformat{\subsection}{\normalfont\normalsize\bfseries}{\thesubsection}{0.5em}{}
\setlength{\parskip}{1pt}
\setlength{\abovedisplayskip}{3pt}
\setlength{\belowdisplayskip}{3pt}
\setlength{\textfloatsep}{5pt}
\setlength{\floatsep}{5pt}
\setlength{\intextsep}{5pt}
\setlength{\columnsep}{14pt}
\pagestyle{empty}
\linespread{0.96}
% let floats fill more of each page so nothing is deferred past the last page
\renewcommand{\topfraction}{0.92}
\renewcommand{\bottomfraction}{0.85}
\renewcommand{\textfraction}{0.06}
\renewcommand{\floatpagefraction}{0.75}
\renewcommand{\dbltopfraction}{0.92}
\renewcommand{\dblfloatpagefraction}{0.72}
\setcounter{topnumber}{4}
\setcounter{bottomnumber}{2}
\setcounter{totalnumber}{6}
\setcounter{dbltopnumber}{3}

\newcommand{\fig}[3]{%
  \begin{figure}[t]\centering\includegraphics[width=#2\linewidth]{figures/#1}%
  \caption{#3}\end{figure}}
\newcommand{\figw}[3]{%
  \begin{figure*}[t]\centering\includegraphics[width=#2\textwidth]{figures/#1}%
  \caption{#3}\end{figure*}}

\begin{document}

\twocolumn[{%
\begin{center}
{\large\bfseries AI4009 Generative AI --- Assignment 1: Latent Spaces and Why GANs Are Hard to Train}
\par\vspace{3pt}
{\normalsize Roll \textbf{25L-8006} \quad$\bullet$\quad seed \textbf{8006} \quad$\bullet$\quad
held-out digit \textbf{6} \quad$\bullet$\quad my six digits \textbf{\{1,\,2,\,5,\,6,\,8,\,9\}}
\quad$\bullet$\quad my attributes \textbf{Smiling + Eyeglasses}}
\par\vspace{8pt}
\end{center}
}]

\section{A --- An AE's latent space is not a space}

\subsection{A1 Train both (CelebA, $d=64$, 20 epochs)}
One encoder--decoder used twice; the models differ only in the reparameterisation, the
two heads and the KL term. Reconstruction is summed over pixels to match the KL summed
over the 64 latent dimensions; a per-pixel mean would shrink it $12288\times$ and the KL
would dominate into posterior collapse. Reported MSE is per pixel.

\begin{center}\small
\begin{tabular}{lcc}
\toprule
& train MSE & val MSE \\
\midrule
AE  & 0.006175 & \textbf{0.006432} \\
VAE & 0.009394 & \textbf{0.009387} \\
\bottomrule
\end{tabular}
\end{center}

The VAE is blurrier by $+0.002955$ MSE/pixel, i.e.\ $1.459\times$ the AE's error, or
$21.92$\,dB vs $20.27$\,dB PSNR --- a gap of $1.64$\,dB. Its KL settled at $57.51$ nats
($\approx$\,83 bits through a 64-d bottleneck) by epoch 5 and held there, so no posterior
collapse. A second difference: the AE ends at
0.006175 train / 0.006432 val with a bouncing validation curve, the VAE at 0.009394 /
0.009387 --- identical and monotone. The KL is also acting as a regulariser.

\fig{A1_recon}{0.88}{\textbf{A1.} Originals over reconstructions, 8 columns, for each model.
The AE keeps the red cap, the background lettering and the green hairband; the VAE
dissolves them.\label{fig:a1}}

\subsection{A2 What the VAE can do: interpolate}
Eleven frames, interpolating the \emph{means} $\mu$ rather than samples from $q$.
\textbf{Every intermediate frame is a plausible face} --- none degenerates into
non-face structure. \textbf{Both attributes change gradually, neither jumps.} For
Smiling the frame-to-frame MSE rises monotonically ($0.000615\rightarrow0.001366$) with
no interior peak. For Eyeglasses it peaks mid-strip at $t=0.4\rightarrow0.5$
($0.002529$), which looks like a jump --- but decomposing by image region shows it is
not: the eye band (rows 20--36) peaks at $t=0.1\rightarrow0.2$ and decays monotonically
thereafter, at $1.36\times$ its median, while the mouth band on the Smiling strip is
flat at $1.06\times$. The mid-strip peak is identity and background morphing, not the
glasses. So the glasses region changes early and concentratedly while the smile changes at a
constant rate; whole-image MSE cannot separate attribute from identity drift.

\figw{A2_vae_interp}{0.80}{\textbf{A2.} VAE latent interpolation between encoded means,
$t=0.0\ldots1.0$. Top: Smiling $+1\rightarrow-1$. Bottom: Eyeglasses
$+1\rightarrow-1$.\label{fig:a2}}

\subsection{A3 What the AE cannot do: generate}
The same 16 codes $z\sim\mathcal{N}(0,I)$ through both decoders (Fig.~\ref{fig:a3}):
the AE returns brown and grey mush, the VAE returns faces. The third panel decodes the
AE's \emph{own} encoder outputs and is the sharpest of the three --- the decoder is not
broken.

The measurement explains it. For a 64-d standard normal $\mathbb{E}\|z\|\approx8$;
measured on 2000 validation images:

\begin{center}\small
\begin{tabular}{lcc}
\toprule
 & mean $\|z\|$ & mean per-dim std \\
\midrule
AE codes        & 32.06 & 3.847 \\
VAE $\mu$       & 6.73  & 0.837 \\
$\mathcal{N}(0,I)$ & 7.97 & 1.000 \\
\bottomrule
\end{tabular}
\end{center}

The AE's codes sit four times further from the origin than any prior draw and the two
histograms in Fig.~\ref{fig:a3n} \textbf{do not overlap at all}: no prior sample lands
in the radial band the decoder was trained on, so it is extrapolating. The VAE's
$\mu$ are slightly under-dispersed (std 0.837), which is correct rather than a fault:
$\mathrm{std}(\mu)^2+\mathbb{E}[\sigma^2]\approx1$ implies $\mathbb{E}[\sigma^2]\approx0.30$,
i.e.\ an encoder genuinely using its noise, consistent with 57.5 nats of KL.

\figw{A3_prior_samples}{0.80}{\textbf{A3.} The same 16 codes through both decoders, and
(right) the AE decoder on codes its own encoder produced.\label{fig:a3}}

\fig{A3_code_norms}{0.98}{\textbf{A3.} Where the codes actually live. The AE cloud
(15--55) and the prior draws (5--11) are disjoint.\label{fig:a3n}}

\section{B --- What the AE is actually good for}

\subsection{B1 Denoising}
Noisy input $\tilde x=\mathrm{clip}(x+0.3\varepsilon,0,1)$, clean target, 10 epochs.
Mean test PSNR: \textbf{13.32\,dB noisy $\rightarrow$ 19.80\,dB denoised}, a gain of
\textbf{6.47\,dB}. Test noise is drawn from a fixed generator so the number is
reproducible across epochs.

\emph{Why can a model that cannot generate a face remove noise from one?} Because
denoising only requires projecting a corrupted input back onto the data manifold, and
the noise is off-manifold, so a narrow bottleneck discards it almost by construction.
Generating requires something strictly stronger --- a distribution over codes that can
be \emph{sampled} --- and A3 shows the AE has no such thing.

\fig{B1_denoise}{0.88}{\textbf{B1.} Clean / noisy / denoised.\label{fig:b1}}

\subsection{B2 The held-out digit as an anomaly}
An AE trained on MNIST with every 6 removed, scored by per-image reconstruction MSE.
ROC-AUC computed from the Mann--Whitney rank statistic (not imported):
\textbf{AUC = 0.883}; mean error 0.00837 on seen digits vs \textbf{0.01846} on the
held-out 6, a ratio of $2.21\times$. The per-digit means (\texttt{figures/B2\_anomaly.png}) put 8 (0.01199) and
2 (0.01159) next, which is the expected confusion given a 6's geometry. Same model,
same bottleneck: useless as a generator, useful as a detector.

\section{C --- How a GAN fails}

\subsection{C1 Watch the gradient die}
Saturating loss $L_G=\log(1-D(G(z)))$, five D steps per G step, $\mathrm{lr}_D=8\mathrm{e}{-4}$
against $\mathrm{lr}_G=2\mathrm{e}{-4}$.

\begin{quote}\small
$\|\nabla_{\theta_G}\| = \mathbf{6.2805\times10^{-5}}$ at iteration 100; it first fell
below 1\% of that ($6.2805\times10^{-7}$) at iteration \textbf{345}.
\end{quote}

A rolling-median (window 51) version of the same rule gives iteration \textbf{761}, so
345 is the first crossing and 761 the sustained one. Both are in
\texttt{logs/C1\_saturating.csv}. Note that the iteration-100 reference is itself
already $135\times$ below the peak of $8.49\times10^{-3}$ at iteration 1 --- measured
from the true peak, the gradient was below 1\% by iteration \textbf{32}. By the end
$D(x)=1.000$, $D(G(z))=9.8\times10^{-11}$ and $\|\nabla_{\theta_G}\|=2.0\times10^{-10}$,
a fall of $4.2\times10^{7}$ from peak. The samples never leave their initialisation
(Fig.~\ref{fig:c2g}, left).

\fig{C1_gradnorm}{0.98}{\textbf{C1.} $D(x)$ and $D(G(z))$, and the generator gradient norm
on a log axis with the death iteration marked.\label{fig:c1}}

\subsection{C2 Fix it with one line}
Identical seed, architecture and D schedule; only $L_G=-\log D(G(z))$.

\begin{center}\small
\begin{tabular}{lcc}
\toprule
median $\|\nabla_{\theta_G}\|$, last 500 it. & C1 & C2 \\
\midrule
 & $6.24\times10^{-10}$ & $3.265$ \\
\bottomrule
\end{tabular}
\end{center}

A ratio of $5.2\times10^{9}$. Both losses share a fixed point, since both are minimised
as $D(G(z))\rightarrow1$, but their gradients differ by a factor of $D(G(z))^2$ in the
regime that matters. The saturating loss carries a prefactor $\propto
D(G(z))/(1-D(G(z)))$, so when a strong D drives $D(G(z))$ to $10^{-10}$ the update
vanishes exactly when the generator is worst and most needs signal. The non-saturating loss scales as $1/D(G(z))$, which \emph{grows} in that regime: under
the identical schedule C2's gradient held at 2--3.4 and rose during training, with
$D(x)=0.894$, $D(G(z))=0.030$ at the end. D still wins; the difference is that G keeps
receiving a usable gradient while losing.

\fig{C2_overlay}{0.80}{\textbf{C2.} Both gradient-norm curves, log axis, same seed and
schedule.\label{fig:c2}}

\figw{C2_grids}{0.64}{\textbf{C2.} The two generators at the same iteration (345/350).
One line of loss separates them.\label{fig:c2g}}

\subsection{C3 Make it collapse, then measure it}
Recipe as given: $z$-dim 2, no BatchNorm in G, five G steps per D step,
$\mathrm{lr}_G=2\mathrm{e}{-3}$ vs $\mathrm{lr}_D=2\mathrm{e}{-4}$. The result is a
\textbf{total} collapse: of 5000 classified samples, \textbf{99.94\% are digit 9}.

\begin{center}\small
\begin{tabular}{lccc}
\toprule
modes covered & reverse-KL & sample std \\
\midrule
\textbf{1 / 6} & \textbf{1.7867} & 0.0129 \\
\bottomrule
\end{tabular}
\end{center}

(Classifier test accuracy 0.9895 on the six digits.) One honest caveat about the stopping
criterion: the samples did \emph{not} stop changing. The MSE between consecutive
fixed-noise grids fell only from 0.926 to 0.804 over 1500 iterations
(\texttt{figures/C3\_movement.png}). Every $z$ maps to essentially the same image at any given
iteration --- that is what the 0.0129 diversity says --- but that single image keeps
drifting from iteration to iteration. The collapse is in class coverage, not in pixel
stasis. $D$ meanwhile sat at chance ($D(x)=0.494$, $D(G(z))=0.464$), which is the
mechanism: G races ahead of a stale D and parks every $z$ on the one sample that fools it.

\fig{C3_samples}{0.50}{\textbf{C3.} 64 samples. One mode, one digit.\label{fig:c3}}

\section{D --- A GAN that actually works}

\subsection{D1 DCGAN}
Strided convolutions, BatchNorm in both networks, LeakyReLU(0.2) in D, ReLU in G, tanh
output, Adam($2\mathrm{e}{-4}$, $\beta_1=0.5$), non-saturating loss, 1:1 steps, 5000
iterations. Result: \textbf{6/6 modes, reverse-KL 0.0182}, histogram
[0.126, 0.203, 0.191, 0.126, 0.163, 0.191] against a uniform target of 0.167. At the end
$D(x)=0.578$, $D(G(z))=0.385$, $d_{\text{loss}}=1.110$ --- a discriminator that genuinely
cannot separate the two distributions.

\figw{D1_progression}{0.72}{\textbf{D1.} The same fixed $z$ at four points in training:
noise, then structure, then digits.\label{fig:d1}}

\subsection{D2 One stabiliser, honestly reported}
Minibatch discrimination (Salimans et al.), written out rather than imported, placed in
the \emph{discriminator}. A unit test confirms it responds to batch similarity: an
identical batch gives a feature of 7.000, a diverse batch 0.228.

\textbf{It did not help.} IS improved marginally (4.206 vs 4.119) but both
distribution-level measures got worse: reverse-KL 0.0665 vs 0.0182, FID 38.75 vs 27.48,
with mode coverage unchanged at 6/6. The reason is in the training logs: the extra
batch-similarity channel strengthened an already dominant discriminator
($D(x)=0.998$ vs $0.578$, $d_{\text{loss}}=0.085$ vs $1.110$, $g_{\text{loss}}=7.0$ vs
$0.86$). D1 had not collapsed, so there was no mode collapse left to fix; the only
effect of the stabiliser was to push the game further from equilibrium.

\figw{D2_histograms}{0.88}{\textbf{D1/D2.} Class histograms for all three models, dashed
line = uniform target.\label{fig:d2}}

\section{E --- Measuring it: IS and FID}

Both metrics implemented from the formulas on the six-digit classifier: its softmax gives
$p(y|x)$, its 128-d penultimate layer gives the FID features. No \texttt{pytorch-fid},
\texttt{clean-fid} or scikit-learn was used at any point, including for cross-checking.
Validation is by cases with known answers: IS of a confident, uniformly distributed
predictor $=K=6$; IS of a maximally unsure predictor $=1$; $\mathrm{FID}(X,X)=0$; and a
mean shift of 3 across 32 dimensions returns 288.2 against an analytic 288.

\begin{center}\small
\begin{tabular}{lcccc}
\toprule
model & modes & rev-KL & IS & FID \\
\midrule
real data      & --    & --     & $5.698\pm0.031$ & 0.32 \\
C1 saturating  & 2/6   & 1.6288 & $1.093\pm0.004$ & 624.31 \\
C3 collapsed   & 1/6   & 1.7867 & $1.042\pm0.002$ & 670.66 \\
D1 DCGAN       & 6/6   & 0.0182 & $4.119\pm0.077$ & \textbf{27.48} \\
D2 + minibatch & 6/6   & 0.0665 & $\mathbf{4.206\pm0.066}$ & 38.75 \\
\bottomrule
\end{tabular}
\end{center}

IS is over 10 splits of 5000 samples; FID is against the real training set (10\,000
images).

\subsection{E2 The matrix square root}
$\Sigma_r\Sigma_g$ is a product of two PSD matrices, so its eigenvalues are real and
non-negative and the true square root is real. An imaginary component can nevertheless
appear from rounding in the Schur decomposition that \texttt{scipy.linalg.sqrtm} uses
internally, which is why one takes \texttt{.real}. In this run the conditioning was
visibly poor --- SciPy raised \texttt{LinAlgWarning: Matrix is singular} on every call,
since a 128-d covariance is estimated from 5000 samples --- but the returned array was
purely real: the recorded imaginary magnitude was \textbf{exactly 0.0} for all five
sets, so \texttt{.real} discarded nothing.

\subsection{E3 Two sanity checks}
\textbf{The floor.} FID between two disjoint halves of the real test set (3000 each) is
\textbf{0.568}, not zero. It measures finite-sample estimation error: $\mu$ and $\Sigma$
computed from 3000 samples differ from those of another 3000 samples drawn from the same
distribution, so 0.568 is the noise level below which no FID difference is meaningful.

\textbf{Monotonicity.} FID against the real set for Gaussian noise
$\sigma=0,0.1,0.25,0.5,1.0$: 0.375, 7.23, 60.14, 265.19, 566.70 --- monotone
(Fig.~\ref{fig:e3}).

\fig{E3_fid_monotonicity}{0.78}{\textbf{E3.} FID rises monotonically with
corruption.\label{fig:e3}}

\subsection{E3(c) What each metric is blind to}
My collapsed model scored IS 1.042 and FID 670.66 against D1's 4.119 and 27.48. IS never
looks at the real data at all --- it only asks whether each sample is confidently
\emph{something} and whether the set spreads across labels --- so it cannot know that
four of my six digits are missing, and it is bounded above by the number of classes
actually covered, which is why a 1-mode generator is capped near 1 regardless of how
clean its 9s are. FID does compare against real data, but only through a mean and a
covariance in a 128-d feature space, so it cannot say \emph{which} digits vanished and
would score two entirely different failures identically if their first two moments
matched. The sharpest evidence is that my C1 model, which emits pure salt-and-pepper
noise, beats the collapsed model on \emph{both} metrics (IS 1.093 vs 1.042, FID 624.31
vs 670.66) --- neither metric can see that one of these generators produces clean,
recognisable digits and the other produces no images at all. Both numbers are also
hundreds of times above my 0.568 floor, so the ranking is not noise; it is what these
statistics genuinely measure.

\section{What surprised me}

\textbf{1. A prior draw is \emph{closer} to the AE's codes than the codes are to each
other.} I expected the opposite. Nearest-neighbour distance from $\mathcal{N}(0,I)$
draws to AE codes is \textbf{16.09}, while the distance between AE codes themselves is
\textbf{25.45} (Fig.~\ref{fig:a3n}). This is not a contradiction of A3, it is 64-d
geometry: the AE's per-dimension mean is $-0.175$, so the codes form a diffuse
\emph{shell} of radius $\approx$\,32 whose points are enormously far apart, and a prior
draw lands near the origin --- the hole in the middle. Being central makes you moderately
close to everything while still sitting where no encoder ever placed a code. Euclidean
nearest-neighbour distance is simply the wrong instrument for ``is this point inside the
data distribution'' when the data lives on a shell; the radial histogram is the right one.

\textbf{2. Pure noise scores better than mode collapse on both metrics.} C1's generator
never left its initialisation and emits salt-and-pepper (Fig.~\ref{fig:c2g}, left), yet
it beats the collapsed C3 model on IS (1.093 vs 1.042) \emph{and} on FID (624.31 vs
670.66). The lecture framing --- IS rewards confident, diverse output, FID punishes
distribution mismatch --- would have predicted C3 ahead on both. What actually happens is
that a diffuse noise cloud overlaps the real feature statistics slightly better than a
tight cluster sitting in the wrong place, and the classifier assigns noise to a couple of
classes with moderate confidence.

\section{What I could not get to work}

\textbf{Isolating the attribute in A2.} Three attempts all measured the identity morph
instead: whole-image frame MSE peaked mid-strip; band-restricted MSE put the eye
region's peak at $t=0.1$--$0.2$; band brightness was dominated by skin and hair and
\emph{fell} across the Eyeglasses strip even as the glasses vanished. The endpoints
differ in every respect, so no global statistic on an image pair separates them. The
right tool is an attribute vector (difference of mean $\mu$ with and without the
attribute, applied to one fixed identity), which A2 did not ask for and I did not run.

\textbf{The C3 stopping criterion.} ``Train until the samples stop changing'' never
became true: consecutive fixed-noise MSE fell only 0.926$\rightarrow$0.804 over 1500
iterations. With $z$-dim 2 and no BatchNorm the generator keeps wandering; what converged
was class coverage, not the images. I report the class collapse as the result and the
movement curve as the caveat rather than claiming a convergence I did not observe.

\textbf{The C1 reference point.} With this D schedule the gradient had already fallen
$135\times$ from its peak by iteration 100, so the handout's threshold describes a late
stage of a collapse mostly over by iteration 32. I report both numbers.

\textbf{The stabiliser.} D2 was worse on two of three measures. I did not tune it
further: with D1 already at 6/6 and reverse-KL 0.018 there was no failure left to fix,
so a better-tuned minibatch discriminator would at best match the baseline.

\section{Declaration of LLM use}

I used Claude (Anthropic) for: scaffolding the Kaggle CelebA loader and the training
loops around the provided starter; explaining the ELBO scaling issue (summed vs
per-pixel reconstruction); debugging an \texttt{IndexError} in my A3 nearest-neighbour
baseline and an \texttt{IndentationError} in the FID function; discussing how to read
the results, including the three failed A2 measurements and the C1-vs-C3 inversion; and
drafting this report's structure from my own numbers. Every run, number and figure here is from my own
execution of my own notebook under seed 8006, and I can derive any part of it.

\end{document}
